% =====================================================================
% Follow-up note to bt_baseline_note v1.0:
% deg_a = 0 row extension and V_quad row-membership re-attribution
% under the extended four-row enumeration.
%
% Standalone SIARC short note (Zenodo cite-target candidate; deposit
% fire is OPERATOR-GATED and is NOT performed by this artefact).
%
% Provenance: this note implements LANE-2 Item 3 verdict
% LEAVE_V1_0_CANONICAL_WITH_VERDICT_AS_FOLLOW_UP_NOTE
% (siarc-relay-bridge dee3c01), with substrate cited from
%   - LANE-2 deposit (bridge dee3c01, six-item verdict + V5 + V6 + P3),
%   - 064 PHASE-A-DEG_A-ZERO-SUPPLEMENTARY (bridge 6a150b6),
%   - 065 PCF2-CF_VALUE-AUDIT-9IMPLS (bridge 6a150b6), and
%   - 066 PCF1-V13-V_QUAD-ROW-REFRAMING (bridge 9261c79).
% No fresh symbolic derivation appears in this note: every claim
% is a verbatim composition of pre-existing substrate.
%
% Date: 2026-05-07
% Filename convention: bt_baseline_note_followup_v1_0
%   (presentation-grade; companion to bt_baseline_note v1.0)
% =====================================================================

\documentclass[11pt,reqno]{amsart}

\usepackage[margin=1in]{geometry}
\usepackage{amsmath,amssymb,amsthm,mathrsfs}
\usepackage{enumitem}
\usepackage{xcolor}
\usepackage{booktabs}
\usepackage{hyperref}
\usepackage[capitalise,noabbrev]{cleveref}

\definecolor{darklink}{rgb}{0.0,0.2,0.5}
\hypersetup{
  colorlinks=true,
  linkcolor=darklink,
  citecolor=darklink,
  urlcolor=darklink,
  pdftitle={Follow-up note to bt_baseline_note v1.0:
            deg_a = 0 row extension and V_quad row-membership
            re-attribution under the extended four-row enumeration},
  pdfauthor={Papanokechi},
  pdfkeywords={bt_baseline_note follow-up, Phase A WZ table,
               deg_a = 0 row, V_quad row-membership, SIARC PCF
               Wallis stratum, mechanism (i') open-content
               closure}
}

\theoremstyle{plain}
\newtheorem{theorem}{Theorem}[section]
\newtheorem{lemma}[theorem]{Lemma}
\newtheorem{proposition}[theorem]{Proposition}
\newtheorem{corollary}[theorem]{Corollary}
\newtheorem{conjecture}[theorem]{Conjecture}

\theoremstyle{definition}
\newtheorem{definition}[theorem]{Definition}
\newtheorem{example}[theorem]{Example}
\newtheorem{problem}[theorem]{Open problem}

\theoremstyle{remark}
\newtheorem{remark}[theorem]{Remark}

\newcommand{\PCF}{\mathrm{PCF}}
\newcommand{\dps}{\mathrm{dps}}
\newcommand{\Anaive}{A_{\mathrm{naive}}}
\newcommand{\Afit}{A_{\mathrm{fit}}}
\newcommand{\dega}{\deg a}
\newcommand{\degb}{\deg b}
\newcommand{\Vquad}{V_{\mathrm{quad}}}

\def\version{1.0 (follow-up)}

\title[Follow-up note to bt\_baseline\_note v1.0]
      {Follow-up note to \texttt{bt\_baseline\_note} v1.0:
       $\dega = 0$ row extension and $\Vquad$ row-membership
       re-attribution under the extended four-row enumeration}

\author{Papanokechi}
\address{Independent researcher, Yokohama, Japan.}

\date{2026-05-07}

\begin{document}

\maketitle

\begin{abstract}
This note is an additive follow-up to \texttt{bt\_baseline\_note}
v1.0 (Zenodo concept DOI
\href{https://doi.org/10.5281/zenodo.20048196}{10.5281/zenodo.20048196};
version DOI
\href{https://doi.org/10.5281/zenodo.20048197}{10.5281/zenodo.20048197}).
The v1.0 main result (\cite[Theorem~1.1]{bt_baseline_note_v10})
is canonical as stated for the band $\dega \in \{d-1, d, d+1\}$
of three SIARC conventions (\cite[\S 2.2 -- \S 2.3]%
{bt_baseline_note_v10}). The present note records two
additive items, both grounded entirely in pre-existing
substrate: (i)~the Phase~A WZ-balance table extends to a fourth
row at $\dega = 0$ (the $(1, b)$ corner inhabited by every
PCF-2 harvest implementation per LANE-2 P1), with
$\Anaive = 2d$ at that corner via Balance~III; and
(ii)~the $\Vquad$ entry of the PCF-1 v1.3 \S 6 Theorem~5 row
table aligns with the $\dega = 0$ row at $d = 2$ as a matter of
row-membership, so the open content of v1.0 \S 6
Q\,[mechanism] (the closure of the borderline-locus posit
\textup{(i$'$)}) admits a row-membership reading. The substrate
is the LANE-2 canonical re-arbitration
(siarc-relay-bridge \cite{lane2_dee3c01}) together with three
implementation reports: 064~\cite{phase_a_deg_a_zero_064},
065~\cite{cf_value_audit_pcf2_9impls_065}, and
066~\cite{pcf1_v13_v_quad_row_reframing_066}. The
v1.0 \texttt{.tex} source is unmodified by this note.
\end{abstract}

\section{Scope and relationship to v1.0}
\label{sec:scope}

This follow-up note is the canonical implementation of the
LANE-2 Item~3 verdict
\texttt{LEAVE\_V1\_0\_CANONICAL\_WITH\_VERDICT\_AS\_FOLLOW\_UP\_NOTE}
(\cite[Item~3 reasoning, L173--200]{lane2_six_item_verdict}).
Per that verdict, \texttt{bt\_baseline\_note} v1.0
Theorem~1.1 is canonical as stated for the band
$\dega \in \{d-1, d, d+1\}$, and the deg\_a = 0 row described
below is an \emph{additive extension}, not a correction. The
v1.0 \texttt{.tex} source (file SHA-256
\verb|6746692C517DC25238473E819527C5682465CDC9E1DEF69D1F6DF31C1014D51B|,
$38{,}023$ bytes; deposited at
\texttt{siarc-relay-bridge/sessions/2026-05-06/T1-PHASE2-BASELINE-NOTE/bt\_baseline\_note.tex})
is unchanged by the present note. The Zenodo v1.0 PDF
(file SHA-256
\verb|23022F0DE77AC8388ED584B2196C0AB995CD8CF18B2DD71EFBC0488A0F6E5B7C|,
$409{,}337$ bytes; concept DOI
\href{https://doi.org/10.5281/zenodo.20048196}{10.5281/zenodo.20048196};
version DOI
\href{https://doi.org/10.5281/zenodo.20048197}{10.5281/zenodo.20048197})
remains the canonical citation target.

The v1.0 \S 4.2 prose at L481--487 of the source file describes
the $\Vquad$ upper-branch case as \emph{consistent with the
borderline-locus mechanism \textup{(i$'$)} of
\texttt{prop:gap-framing}, the closure of which is the open
content of \S 6 Q\,[mechanism]} (\cite[L481--487]%
{bt_baseline_note_v10}). That prose is forward-looking
open-conjecture content as v1.0 itself flags it; the present
note resolves the open content of \S 6 Q\,[mechanism] via a
row-membership reading under the extended four-row Phase~A
enumeration of relay~064~\cite{phase_a_deg_a_zero_064}, with
no edit to the v1.0 source.

\section{Phase A WZ-table extension to $\dega = 0$}
\label{sec:phase-a-extension}

The original Phase~A WZ-balance enumeration of v1.0 covers the
three SIARC conventions $\alpha$ / symmetric / $\delta$ at
$\dega \in \{d-1, d, d+1\}$, $d \in [2, 8]$
(\cite[\S 2.3, \S 3.4]{bt_baseline_note_v10}). Relay~064
records the WZ Newton-polygon balance at the corner case
$\dega = 0$ -- the $(1, b)$ continued-fraction corner used by
every PCF-2 harvest implementation per LANE-2 P1
(\cite[\S 2.2, \S 4]{phase_a_deg_a_zero_064}). The extension
is a single new row at each $d$ and is recorded verbatim in
064 \S 2.3 as the bold rows of the twelve-row enumeration
(file SHA-256
\verb|80E28568FF142B1A81C7C443368349A07A8235CE122D6C6140F88AA0D0706097|,
$16{,}792$ bytes).

The four-row enumeration $\dega \in \{0, d-1, d, d+1\}$ at
each $d$ takes the form (reproduced verbatim from
\cite[\S 2.3, L102--118]{phase_a_deg_a_zero_064}):

\begin{table}[h]
\centering
\small
\begin{tabular}{cccccccc}
\toprule
$d$ & Convention & $\dega$ & $\degb$
    & $\mu_{\mathrm{dom}}$ & $\mu_{\mathrm{sub}}$ & $\Anaive$ \\
\midrule
$2$ & $(1, b)$ corner & $0$ & $2$ & $2$ & $-2$ & $\mathbf{4}$ \\
$2$ & $\alpha$        & $1$ & $2$ & $2$ & $-1$ & $3$ \\
$2$ & symmetric       & $2$ & $2$ & $2$ &  $0$ & $2$ \\
$2$ & $\delta$        & $3$ & $2$ & $2$ &  $1$ & $1$ \\
\midrule
$3$ & $(1, b)$ corner & $0$ & $3$ & $3$ & $-3$ & $\mathbf{6}$ \\
$3$ & $\alpha$        & $2$ & $3$ & $3$ & $-1$ & $4$ \\
$3$ & symmetric       & $3$ & $3$ & $3$ &  $0$ & $3$ \\
$3$ & $\delta$        & $4$ & $3$ & $3$ &  $1$ & $2$ \\
\midrule
$4$ & $(1, b)$ corner & $0$ & $4$ & $4$ & $-4$ & $\mathbf{8}$ \\
$4$ & $\alpha$        & $3$ & $4$ & $4$ & $-1$ & $5$ \\
$4$ & symmetric       & $4$ & $4$ & $4$ &  $0$ & $4$ \\
$4$ & $\delta$        & $5$ & $4$ & $4$ &  $1$ & $3$ \\
\bottomrule
\end{tabular}
\caption{The extended four-row Phase~A WZ-balance enumeration
at $d \in \{2, 3, 4\}$, reproduced verbatim from
\cite[\S 2.3, L102--118]{phase_a_deg_a_zero_064}. Bold rows are
the new $\dega = 0$ rows added by relay~064; the other rows
are reproduced verbatim from the original Phase~A summary at
\texttt{phase\_a\_summary.md} L34--44 via LANE-2 P2 substrate
\cite[L77--87]{lane2_independent_depth_probe}.}
\label{tab:fourrow}
\end{table}

The new bold rows are derived in LANE-2 V6
(\cite[\S V6, L210--308]{lane2_independent_substrate_verification},
file SHA-256
\verb|56063BF7BA8AD6A01A89FA30A3C61FCE68F4AAB0BA0DC7C8E561A81188D7B1F5|,
$15{,}695$ bytes) as Balance~III applied at the $(1, b)$ corner
($c_a = 1$, $d_a = 0$). The relevant V6 conclusion is quoted
verbatim by 064 \S 3 (\cite[\S 3]{phase_a_deg_a_zero_064}):

\begin{quote}
``$\Anaive = \mu_{\mathrm{dom}} - \mu_{\mathrm{sub}} = d - (-d) =
2d$ (when $\dega = 0$). \dots General formula:
$\Anaive = 2d - d_a$.''
\hfill (\cite[V6 Step~4, L274--282]{lane2_independent_substrate_verification})
\end{quote}

\noindent
For $d = 2$, $\dega = 0$: $\Anaive = 4$. For $d = 3$:
$\Anaive = 6$. For $d = 4$: $\Anaive = 8$. The sign of the
recessive root in the same balance is recorded verbatim in V6
Step~3 at L260--266 of the same file as
$\gamma_{\mathrm{sub}} \propto -c_a / c_b$, with the corner
specialisation $\gamma_{\mathrm{sub}} = -1/c_b \cdot e^{d}$ at
$c_a = 1, d_a = 0$.

\section{$\Vquad$ as $\dega = 0$ row member at $d = 2$}
\label{sec:vquad-row}

The $\Vquad$ representative of the PCF-1 v1.3 \S 6 Theorem~5
row table (\cite{papanokechi_pcf1_v13}) carries the empirical
exponent $A = 4$, with $\alpha$-prediction match to
$13$ digits at $\dps = 2200$
(\cite[L566--577]{papanokechi_pcf1_v13}, reproduced verbatim
in \cite[\S 2.2]{pcf1_v13_v_quad_row_reframing_066}). The
$\Vquad$ coefficient declaration in the SIARC audit script
\texttt{algebraic\_independence\_audit.py} L37--40 reads
\verb|VQUAD_ALPHA = [1]| (\cite[V5, L161--179]%
{lane2_independent_substrate_verification}; reproduced in
\cite[\S 3]{pcf1_v13_v_quad_row_reframing_066}), so $a_n = 1$
for all $n$ at the source-code declaration level, hence
$\dega = 0$.

LANE-2 P3 records the substrate-level identification: the
$\Vquad$ family has $a_n \equiv 1$ at the source-code level
(\cite[L77--96]{lane2_independent_depth_probe}). Combined
with~\Cref{tab:fourrow}, the empirical $A = 4$ entry of
$\Vquad$ at $d = 2$ aligns row-by-row with the $\dega = 0$
row of the four-row Phase~A enumeration. This row-membership
re-attribution is the canonical content of relay~066
(\cite[\S 5]{pcf1_v13_v_quad_row_reframing_066}; file SHA-256
\verb|79933B694DD2BF99793429B1A122A4BFF3260A42A93680886D5EF89B4E10FDCD|,
$24{,}073$ bytes).

The framing throughout is \emph{row-membership re-attribution
under extended enumeration}, additive to v1.0. The PCF-1 v1.3
\S 6 Theorem~5 row-table source remains canonical; an explicit
PCF-1 v1.4 amendment is forward-pointed under the
\texttt{PCF1-V13-RECONCILE} jurisdiction
(\cite[\S 1, \S 6]{pcf1_v13_v_quad_row_reframing_066}) and is
not addressed by the present note.

\section{PCF-2 cf\_value implementation uniformity audit}
\label{sec:pcf2-uniformity}

Relay~065 audits 13 cf\_value-class evaluation impls across the
\texttt{pcf-research/pcf2/} layer (a strictly stronger coverage
than the LANE-2 P1 enumeration of nine impls
\cite[L37--74]{lane2_independent_depth_probe}). The aggregate
verdict is \texttt{UNIFORM} at the effective-use layer:
$12$ of $13$ impls classify as \texttt{HC1} (strict inline
\verb|mp.mpf(1)| for $a_n$); $1$ of $13$ classifies as
\texttt{HC0} (\texttt{evaluate\_cf} parameterised signature with
default $\lambda$ and zero non-default call sites);
$0$ of $13$ classifies as \texttt{PARAM}
(\cite[\S 4]{cf_value_audit_pcf2_9impls_065}; file SHA-256
\verb|16512BCC71C9A19ED59B808D8D070877F2160C5A7A062381C0FC9533182D3BF9|,
$11{,}700$ bytes). At the implementation layer of every
PCF-2 R1.1 + R1.3 + Q1 harvest pipeline, the $\dega = 0$
hardcoding is uniform; this is consistent with the $\Vquad$
identification of~\Cref{sec:vquad-row} and supports a
uniform $\dega = 0$ row-membership reading at the $d = 3$ and
$d = 4$ strata as well, alongside the empirical
$\Afit \approx 2d$ records at those strata
(\cite{papanokechi_pcf2_v13} R1.1 + R1.3 + Q1; reproduced in
\cite[\S 4.1]{bt_baseline_note_v10}).

\section{Mechanism \textup{(i$'$)} open-content closure under
the $\dega = 0$ row}
\label{sec:closure}

\paragraph{The v1.0 open-content sentence (verbatim).}
The open content of v1.0 \S 4.2 reads, at L481--487 of
\texttt{bt\_baseline\_note.tex}
(\cite[L481--487]{bt_baseline_note_v10}):

\begin{quote}
``The $\Vquad$ upper branch is consistent with the
borderline-locus mechanism \textup{(i$'$)} of
\texttt{prop:gap-framing} -- the closure of which is the open
content of \S 6 Q\,[mechanism].''
\end{quote}

\noindent
The v1.0 source phrases this as conditional, forward-looking
open-content language; the closure of mechanism~\textup{(i$'$)}
on the SIARC stratum is flagged in v1.0 \S 6 as an open
problem (\texttt{q:mechanism}; \cite[\S 6]{bt_baseline_note_v10}).

\paragraph{Closure under the $\dega = 0$ row reading.}
Under the extended four-row enumeration
of~\Cref{tab:fourrow}, the $\Vquad$ entry $A = 4$ at $d = 2$
aligns with the $\dega = 0$ row of the Phase~A table, with
$A_{\mathrm{naive}} = 2d$ via Balance~III at the corner
$c_a = 1$, $d_a = 0$ (\Cref{sec:phase-a-extension}). This
row-corner placement is independent of any borderline-locus
posit: the $\dega = 0$ row is a normal-case Balance~III row,
and Balance~III is one of the three explicit balance branches
of the v1.0 Wimp--Zeilberger enumeration
(\cite[\S 3.2]{bt_baseline_note_v10}).
A concrete consequence: under the row-membership reading, the
v1.0 \S 6 open problem \texttt{q:mechanism} admits a closure
in which $\Vquad$'s upper branch is realised as a row-corner
of the extended Phase~A enumeration, rather than as a
borderline-locus phenomenon distinct from row-membership. The
borderline-locus posit \textup{(i$'$)}, considered as a
candidate structural mechanism on \emph{other} SIARC strata
where the $\dega = 0$ row reading does not apply, remains an
open future inquiry; the present note does not address it.

\paragraph{Naming the gap.}
The descriptive term ``protocol-to-stratum mismatch'' (the gap
between the protocol-declared $\dega$-domain of PCF-2 v1.3
\S 6 -- $\dega \in \{0, 1\}$ -- and the actually-evaluated
$\dega$ at the implementation layer, uniformly $\dega = 0$ per
\Cref{sec:pcf2-uniformity}) is a candidate label for the gap
that the row-membership reading addresses;\footnote{Term
pending W21 canonical T1-Synth (LANE-1) rule5-vocabulary
ratification per LANE-2 Item~4 verdict
\texttt{DEFER\_TO\_W21}
(\cite[Item~4, L210--213]{lane2_six_item_verdict}). The
present note uses the term descriptively only.} the substrate
basis for the term is LANE-2 V3 + V4 + V5 + P4
(\cite[\S V3 -- \S V5]{lane2_independent_substrate_verification},
\cite[\S P4]{lane2_independent_depth_probe}).

\section{Limitations and open questions}
\label{sec:limitations}

The row-membership reading of~\Cref{sec:vquad-row} and~%
\Cref{sec:closure} carries the following deliberate scope
boundaries. None of the items below is closed by the present
note.

\begin{itemize}[leftmargin=2.2em]
\item The $\dega = 0$ row reading identifies the formal-level
  baseline $\Anaive = 2d$ at the $\dega = 0$ corner under
  Phase~A enumeration. The empirical record carried by
  \cite{papanokechi_pcf2_v13} R1.1 + R1.3 + Q1 records
  $\Afit \approx 2d$ at $d = 3, 4$ at the analytic level
  (\cite[\S 4.1]{bt_baseline_note_v10}). The
  formal-to-analytic sectorial-summability upgrade
  (Borel--Laplace radius identification at the SIARC stratum,
  \cite[ch.~5]{costin2008asymptotics_followup}) is the
  jurisdiction of T1 Phase~3 (LANE-2 Item~2 sub-task 3-D
  authorisation \cite[Item~2, sub-task 3-D]%
  {lane2_six_item_verdict}; gated on a separate relay).
\item The borderline-locus posit \textup{(i$'$)} as a
  structural mechanism on SIARC strata distinct from the
  $(1, b)$ corner remains an open future inquiry. The present
  note records the row-membership reading as a closure of the
  v1.0 \S 6 \texttt{q:mechanism} content for the $\Vquad$
  $d = 2$ case under the extended enumeration; it does not
  rule out structural roles for \textup{(i$'$)} elsewhere.
\item An explicit PCF-1 v1.4 \S 6 row-table amendment that
  renders the $\dega = 0$ row reading at the source-prose
  level is forward-pointed under the
  \texttt{PCF1-V13-RECONCILE} jurisdiction
  (\cite[\S 6]{pcf1_v13_v_quad_row_reframing_066}). The
  PCF-1 v1.3 source \texttt{p12\_pcf1\_main.tex} (file
  SHA-256
  \verb|E83BB377F297DBF0837FACBA257F227DF4579E6A3518C139E3146F17174BE301|)
  remains canonical and is unmodified by this note.
\end{itemize}

\section*{Acknowledgements}

The author thanks the LANE-2 reviewer (Copilot Researcher
Agent) for the canonical re-arbitration recorded at
siarc-relay-bridge \cite{lane2_dee3c01}; and acknowledges the
implementation-layer reports 064~\cite{phase_a_deg_a_zero_064},
065~\cite{cf_value_audit_pcf2_9impls_065}, and
066~\cite{pcf1_v13_v_quad_row_reframing_066} as the substrate
of the present additive follow-up. Every non-citation
mathematical claim in this note is a verbatim composition of
substrate already deposited in those four reports.

\section*{AI Disclosure}

During the preparation of this work the author used GitHub
Copilot (Microsoft) and Anthropic Claude to assist with
manuscript drafting. After using these tools, the author
reviewed and edited the content as needed and takes full
responsibility for the content.

\bibliographystyle{plain}
\bibliography{annotated_bibliography_followup}

\end{document}
